% Generated draft: requires mathematical and source review.
\documentclass[11pt]{article}
\usepackage[margin=25mm]{geometry}
\usepackage{fontspec}
\setmainfont{Latin Modern Roman}
\usepackage{amsmath,amssymb,mathtools,unicode-math}
\setmathfont{Latin Modern Math}
\usepackage{enumitem,needspace,xurl,hyperref}
\setlength{\parindent}{0pt}
\setlength{\parskip}{5pt}
\setlength{\emergencystretch}{3em}
\newcommand{\N}{\mathbb{N}}
\newcommand{\Z}{\mathbb{Z}}
\newcommand{\Q}{\mathbb{Q}}
\newcommand{\R}{\mathbb{R}}
\newcommand{\C}{\mathbb{C}}
\newcommand{\F}{\mathbb{F}}
\newcommand{\A}{\mathbb{A}}
\newcommand{\PP}{\mathbb{P}}
\newcommand{\E}{\mathbb{E}}
\newcommand{\eps}{\varepsilon}
\newcommand{\dd}{\,\mathrm{d}}
\begin{document}
\section*{UnsolvedMath ID: 30000855}
\textbf{Source problem code:} catalogue entry
\subsection{Definitions and mathematical statement}
The supplied question concerns surfaces \(S\) satisfying
\[
K_S^2<3\chi(S)
\]
and having no irregular \(\acute{e}tale\) cover. In the customary smooth projective complex surface setting, \(K_S\) is the canonical divisor,
\[
\chi(S)=\chi(S,\mathcal O_S)=1-h^1(S,\mathcal O_S)+h^2(S,\mathcal O_S),
\qquad q(S)=h^1(S,\mathcal O_S),
\]
and the absence of an irregular \(\acute{e}tale\) cover means that every connected finite \(\acute{e}tale\) cover \(S'\to S\) has \(q(S')=0\). Let \(\pi_1^{\mathrm{alg}}(S)\) denote the profinite \(\acute{e}tale\) fundamental group classifying connected finite \(\acute{e}tale\) covers.

The question asks whether, up to finitely many exceptions, every surface in the intended class satisfies
\[
\pi_1^{\mathrm{alg}}(S)\hookrightarrow(\Z/2\Z)^2.
\]
For profinite groups this asks that the group be isomorphic to one of \(1\), \(\Z/2\Z\), and \((\Z/2\Z)^2\), with the finite target given its discrete topology.

The source extraction does not explicitly state smoothness, projectivity, ground field, minimality, or general type. The conventional definitions above explain the likely invariant notation without asserting that these supply the full intended surface class. Most importantly, ``finitely many exceptions'' is not quantified: it could refer to surfaces up to isomorphism, deformation families, or numerical invariant pairs. These are different claims, so no particular exceptional parameter space is selected. A precise version requires a specified class and equivalence relation for which a finite exceptional set is then allowed.
\subsection{Short English statement}
Apart from a finite exceptional set in an unspecified sense, must a surface with \(K_S^2<3\chi(S)\) and no irregular finite \(\acute{e}tale\) cover have algebraic fundamental group trivial, cyclic of order two, or elementary abelian of order four?
\subsection{Sources}
\begin{itemize}
\item Local UnsolvedMath registry, permanent ID \(30000855\), source code catalogue entry, ``Algebraic Fundamental Groups Below the Noether Line.'' No individual proposer is supplied.
\item \emph{Komplexe Algebraische Geometrie} (2008), Oberwolfach Reports, DOI 10.4171/owr/2007/47, pp. 2791--2850: \url{https://doi.org/10.4171/owr/2007/47}.
\item Supplied reference, ``Mendes Lopes--Pardini, algebraic fundamental groups'': \url{https://arxiv.org/abs/math/0512483}.
\item Supplied literature triage, checked 2026-08-21: status ``open.'' It calls the proposed subgroup description a stronger form of an unresolved surface-group problem and says lower-slope boundedness results do not establish it. This dated assessment derives from supplied OWR curation; time-sensitive status and cited claims require independent verification.
\end{itemize}
\end{document}
